% Chapter 7: Matched-Pairs Experiment
% A First Course in Causal Inference - Peng Ding
% Draft V1 - Stein-style rewrite

\chapter{配对随机实验}\label{ch:7}

\section{引言：从分层到配对的极致追求}\label{sec:7-intro}

在第 \ref{ch:5} 章的分层随机实验（SRE）中，我们将单元按协变量分层，在每层内独立进行完全随机实验。当每层只有一个治疗单元和一个对照单元时，我们得到了 SRE 的最极端形式——\textbf{配对随机实验（Matched-Pairs Experiment, MPE）}。此时，每一层也称为一个\textbf{对（pair）}。

MPE 看似只是 SRE 的一个特例，但它拥有独特估计和推断策略，且与观测性研究中的匹配方法密切相关（将在第 \ref{ch:15} 章详细讨论）。因此，本书将其独立成章。

\textbf{核心问题}：当协变量能够预测潜在结果时，将相似的单元配对后再随机分配治疗，能否进一步提高估计精度？

\section{MPE 的定义与记号}\label{sec:7-definition}

考虑一个包含 $2n$ 个单元的实验。如果我们拥有能够预测结果的协变量，可以根据协变量的相似性将单元配对。对于单标量协变量，我们可以根据该协变量对单元排序，然后配对相邻的单元。对于多维协变量，可以定义单元之间的成对距离，然后使用贪婪算法或最优非二分匹配算法进行配对。\footnote{配对算法的详细内容见 \cite{Greevy:2004}。}

记 $(i,j)$ 表示第 $i$ 对中的第 $j$ 个单元，其中 $i = 1, \dots, n$，$j = 1, 2$。单元 $(i,j)$ 的潜在结果为 $Y_{ij}(1)$（接受治疗）和 $Y_{ij}(0)$（接受对照）。在每一对内，随机分配一个单元接受治疗，另一个接受对照。定义
\begin{equation}
Z_i =
\begin{cases}
1, & \text{若第一个单元接受治疗}, \\
0, & \text{若第二个单元接受治疗}.
\end{cases}
\label{eq:7-treatment-indicator}
\end{equation}

\begin{Definition}[配对随机实验 (MPE)]\label{def:7-mpe}
我们有
\[
(Z_i)_{i=1}^n \stackrel{\text{IID}}{\sim} \text{Bernoulli}(1/2).
\]
即，治疗指示变量 $(Z_1, \dots, Z_n)$ 独立同分布，服从参数为 $1/2$ 的伯努利分布。
\end{Definition}

第 $i$ 对内的观测结果为
\[
Y_{i1} = Z_i Y_{i1}(1) + (1 - Z_i) Y_{i1}(0), \quad
Y_{i2} = Z_i Y_{i2}(0) + (1 - Z_i) Y_{i2}(1).
\]
因此，观测数据为 $(Z_i, Y_{i1}, Y_{i2})_{i=1}^n$。

\section{Fisher 随机化检验}\label{sec:7-frt}

与前述实验一样，我们始终可以使用 Fisher 随机化检验（FRT）来检验\textbf{尖锐零假设}：
\begin{equation}
H_{0F}: Y_{ij}(1) = Y_{ij}(0) \quad \text{对所有 } i = 1, \dots, n \text{ 和 } j = 1, 2.
\label{eq:7-sharp-null}
\end{equation}

在 FRT 中，我们需要从 \cref{eq:7-treatment-indicator} 模拟 $(Z_1, \ldots, Z_n)$ 的分布。以下是基于配对内治疗组与对照组结果之差的几种经典检验统计量。

\subsection{配对内差异}\label{sec:7-within-pair-diff}

定义配对 $i$ 内的\textbf{配对差异}：
\begin{equation}
\hat\tau_i = \text{治疗组结果} - \text{对照组结果} = (2Z_i - 1)(Y_{i1} - Y_{i2}) = S_i(Y_{i1} - Y_{i2}),
\label{eq:7-pair-diff}
\end{equation}
其中 $S_i = 2Z_i - 1$ 是独立同分布的随机符号，均值为 $0$，方差为 $1$。由于 $\hat\tau_i = 0$ 的配对对随机化分布没有贡献，我们在 FRT 讨论中剔除这些配对。

配对差异的均值为
\begin{equation}
\hat\tau = n^{-1} \sum_{i=1}^n \hat\tau_i.
\label{eq:7-avg-tau}
\end{equation}

\subsection{配对 $t$ 统计量}\label{sec:7-paired-t}

\begin{Example}[配对 $t$ 统计量]\label{ex:7-paired-t}
在 $H_{0F}$ 下，$\hat\tau_i$ 是固定的，而 $S_i$ 是随机的。因此
\[
\mathbb{E}(\hat\tau) = 0, \quad
\text{var}(\hat\tau) = n^{-2} \sum_{i=1}^n \text{var}(S_i)(Y_{i1} - Y_{i2})^2 = n^{-2} \sum_{i=1}^n \hat\tau_i^2.
\]

基于独立随机变量之和的中心极限定理，我们有正态近似：\footnote{推导见附录 \cref{sec:appendix-7-clt}}
\begin{equation}
\frac{\hat\tau}{\sqrt{n^{-2} \sum_{i=1}^n \hat\tau_i^2}} \to \mathrm{N}(0, 1)
\label{eq:7-t-asymptotic}
\end{equation}
依分布收敛。许多标准教科书建议在 MPE 中使用以下配对 $t$ 统计量：
\begin{equation}
t_{\text{pair}} = \frac{\hat\tau}{\sqrt{\{n(n-1)\}^{-1} \sum_{i=1}^n (\hat\tau_i - \hat\tau)^2}}.
\label{eq:7-t-pair}
\end{equation}
在 $H_{0F}$ 下，当 $n$ 较大且 $\hat\tau$ 较小时，$t_{\text{pair}}$ 与 \cref{eq:7-t-asymptotic} 几乎相同。

在经典统计学中，使用 $t_{\text{pair}}$ 的动机基于不同框架。当 $\hat\tau_i \stackrel{\text{IID}}{\sim} \mathrm{N}(0, \sigma^2)$ 时，可以证明 $t_{\text{pair}} \sim t(n-1)$，即 $t_{\text{pair}}$ 的精确分布是自由度为 $n-1$ 的 $t$ 分布，在大 $n$ 下接近 $\mathrm{N}(0, 1)$。R 函数 \texttt{t.test} 可以实现此检验。\cref{ex:7-paired-t} 的讨论给出了经典配对 $t$ 检验的另一种证明，无需假设数据的正态性。
\end{Example}

\subsection{Wilcoxon 符号秩统计量}\label{sec:7-wilcoxon}

\begin{Example}[Wilcoxon 符号秩统计量]\label{ex:7-wilcoxon}
基于 $(|\hat\tau_1|, \dots, |\hat\tau_n|)$ 的秩 $(R_1, \ldots, R_n)$，我们可以定义检验统计量
\begin{equation}
W = \sum_{i=1}^n \mathbb{I}(\hat\tau_i > 0) R_i.
\label{eq:7-wilcoxon}
\end{equation}
在 $H_{0F}$ 下，$|\hat\tau_i|$ 是固定的，因此 $R_i$ 也是固定的。由此可得
\[
\mathbb{E}(W) = \frac{1}{2} \sum_{i=1}^n R_i = \frac{1}{2} \sum_{i=1}^n i = \frac{n(n+1)}{4},
\]
\[
\text{var}(W) = \frac{1}{4} \sum_{i=1}^n R_i^2 = \frac{1}{4} \sum_{i=1}^n i^2 = \frac{n(n+1)(2n+1)}{24}.
\]

独立随机变量之和的中心极限定理保证以下正态近似：\footnote{推导见附录 \cref{sec:appendix-7-wilcoxon}}
\begin{equation}
\frac{W - n(n+1)/4}{\sqrt{n(n+1)(2n+1)/24}} \to \mathrm{N}(0, 1)
\label{eq:7-wilcoxon-asymptotic}
\end{equation}
依分布收敛。R 函数 \texttt{wilcox.test} 可以实现精确和渐近检验。
\end{Example}

\subsection{Kolmogorov--Smirnov 型统计量}\label{sec:7-ks}

\begin{Example}[Kolmogorov--Smirnov 型统计量]\label{ex:7-ks}
在 $H_{0F}$ 下，绝对值 $(|\hat\tau_1|, \dots, |\hat\tau_n|)$ 是固定的，但符号是随机的。因此 $(\hat\tau_1, \dots, \hat\tau_n)$ 和 $-(\hat\tau_1, \dots, \hat\tau_n)$ 应具有相同分布。

设
\[
\hat F(t) = n^{-1} \sum_{i=1}^n \mathbb{I}(\hat\tau_i \leq t)
\]
是 $(\hat\tau_1, \dots, \hat\tau_n)$ 的经验分布函数，
\[
1 - \hat F(-t-) = n^{-1} \sum_{i=1}^n \mathbb{I}(-\hat\tau_i \leq t)
\]
是 $-(\hat\tau_1, \dots, \hat\tau_n)$ 的经验分布函数，其中 $\hat F(-t-)$ 是 $\hat F(\cdot)$ 在 $-t$ 处的左极限。Kolmogorov--Smirnov 型统计量定义为
\begin{equation}
D = \max_t |\hat F(t) + \hat F(-t-) - 1|.
\label{eq:7-ks-statistic}
\end{equation}
\citet{Butler:1969} 提出了这个检验统计量并推导了其精确和渐近分布。标准软件包中未实现此统计量，但我们可以基于 FRT 模拟其精确分布并计算 p 值。
\end{Example}

\subsection{符号统计量}\label{sec:7-sign}

\begin{Example}[符号统计量]\label{ex:7-sign}
符号统计量仅使用配对内差异的符号：
\begin{equation}
\Delta = \sum_{i=1}^n \mathbb{I}(\hat\tau_i > 0).
\label{eq:7-sign-stat}
\end{equation}
在 $H_{0F}$ 下，
\[
\mathbb{I}(\hat\tau_i > 0) \stackrel{\text{IID}}{\sim} \text{Bernoulli}(1/2),
\]
因此 $\Delta \sim \text{Binomial}(n, 1/2)$。这给出了一个基于二项分布的精确检验，R 函数 \texttt{binom.test} 可以实现。利用中心极限定理，我们还可以基于以下正态近似进行检验：\footnote{推导见附录 \cref{sec:appendix-7-sign}}
\begin{equation}
\frac{\Delta - n/2}{\sqrt{n/4}} \to \mathrm{N}(0, 1)
\label{eq:7-sign-asymptotic}
\end{equation}
依分布收敛。
\end{Example}

\subsection{二值结果的 McNemar 统计量}\label{sec:7-mcnemar}

当结果为二值时，我们可以更紧凑地汇总 MPE 的观测数据。给定一对，治疗结果和对照结果可以各为 $1$ 或 $0$，形成如下 $2 \times 2$ 列联表：

\begin{center}
\begin{tabular}{l|cc}
 & \text{对照结果}=1 & \text{对照结果}=0 \\ \hline
\text{治疗结果}=1 & $m_{11}$ & $m_{10}$ \\
\text{治疗结果}=0 & $m_{01}$ & $m_{00}$
\end{tabular}
\end{center}

\begin{Example}[McNemar统计量]\label{ex:7-mcnemar}
在 $H_{0F}$ 下，一致对的数量 $m_{11}$ 和 $m_{00}$ 是固定的，$m_{10} + m_{01}$ 也是固定的。因此唯一的随机分量是 $m_{10}$，其分布为
\[
m_{10} \sim \text{Binomial}(m_{10} + m_{01}, 1/2).
\]

这给出了一个基于二项分布的精确检验。R 函数 \texttt{mcnemar.test} 给出基于二项分布正态近似的渐近检验：
\begin{equation}
\frac{m_{10} - (m_{10} + m_{01})/2}{\sqrt{(m_{10} + m_{01})/4}} = \frac{m_{10} - m_{01}}{\sqrt{m_{10} + m_{01}}} \to \mathrm{N}(0, 1)
\label{eq:7-mcnemar}
\end{equation}
依分布收敛。精确 FRT 和渐近检验都不依赖于 $m_{11}$ 或 $m_{00}$，只有不一致对的数量在这些检验中起作用。
\end{Example}

\section{Neymanian 推断}\label{sec:7-neymanian}

第 $i$ 对内的平均因果效应为
\begin{equation}
\tau_i = \frac{1}{2}\{Y_{i1}(1) + Y_{i2}(1) - Y_{i1}(0) - Y_{i2}(0)\},
\label{eq:7-pair-ace}
\end{equation}
所有单元的平均因果效应为
\begin{equation}
\tau = n^{-1} \sum_{i=1}^n \tau_i = (2n)^{-1} \sum_{i=1}^n \sum_{j=1}^2 \{Y_{ij}(1) - Y_{ij}(0)\}.
\label{eq:7-total-ace}
\end{equation}

直观上，$\hat\tau_i$ 是 $\tau_i$ 的无偏估计，因此 $\hat\tau$ 是 $\tau$ 的无偏估计。\footnote{完整推导见附录 \cref{sec:appendix-7-variance}}

然而，在 MPE 下，我们无法遵循 SRE 的策略来估计 $\hat\tau$ 的方差。因为在每一对内，我们只有一个治疗单元和一个对照单元，结果的配对内样本方差无法定义。数据不允许我们在配对 $i$ 内估计 $\hat\tau_i$ 的方差。

能否在 MPE 中估计 $\hat\tau$ 的方差？让我们换一个视角，忘记 MPE，转向经典独立同分布抽样。如果 $\hat\tau_i$ 是均值为 $\mu$、方差为 $\sigma^2$ 的独立同分布变量，那么 $\hat\tau = n^{-1}\sum_{i=1}^n \hat\tau_i$ 的方差是 $\sigma^2/n$。$\sigma^2$ 的无偏估计是样本方差 $(n-1)^{-1}\sum_{i=1}^n (\hat\tau_i - \hat\tau)^2$，因此 $\text{var}(\hat\tau)$ 的无偏估计是
\begin{equation}
\hat V = \{n(n-1)\}^{-1} \sum_{i=1}^n (\hat\tau_i - \hat\tau)^2.
\label{eq:7-variance-estimator}
\end{equation}
上述讨论似乎与 MPE 完全不同。但至少它启发我们使用 $\hat V$，即用 $\hat\tau_i$ 的配对间方差来估计 $\hat\tau$ 的方差。当然，它是在不同假设下推导的。这在 MPE 中也适用吗？\cref{th:7-variance-conservative} 给出了肯定的回答。

\begin{Theorem}[方差估计的保守性]\label{th:7-variance-conservative}
在 MPE 下，由 \cref{eq:7-variance-estimator} 定义的 $\hat V$ 是 $\hat\tau$ 真实方差的一个保守估计：
\begin{equation}
\mathbb{E}(\hat V) - \text{var}(\hat\tau) = \{n(n-1)\}^{-1} \sum_{i=1}^n (\tau_i - \tau)^2 \geq 0.
\label{eq:7-variance-bias}
\end{equation}
当 $\tau_i$ 在各配对间恒定时，$\mathbb{E}(\hat V) = \text{var}(\hat\tau)$。
\end{Theorem}

\cref{th:7-variance-conservative} 表明，在 MPE 下，$\hat V$ 通常是保守的方差估计量，只有当各配对的平均因果效应恒定时才是无偏的。这个结果有些令人惊讶，因为 $\hat V$ 依赖于 $\hat\tau_i$ 的配对间方差，而 $\text{var}(\hat\tau)$ 依赖于每个 $\hat\tau_i$ 的配对内方差。\footnote{证明见附录 \cref{sec:appendix-7-theorem-proof}}

与前述实验一样，Neymanian 方法依赖大样本近似：
\begin{equation}
\frac{\hat\tau - \tau}{\sqrt{\text{var}(\hat\tau)}} \to \mathrm{N}(0, 1)
\label{eq:7-clt}
\end{equation}
在 $n \to \infty$ 且某些正则条件成立时依分布收敛。由于方差的过度估计，Wald 型置信区间
\begin{equation}
\hat\tau \pm z_{1-\alpha/2} \sqrt{\hat V}
\label{eq:7-wald-ci}
\end{equation}
以至少 $1 - \alpha$ 的概率覆盖 $\tau$。

点估计量 $\hat\tau$ 和方差估计量 $\hat V$ 都可以方便地通过 OLS 得到，如下命题所示。

\begin{Proposition}[OLS 解释]\label{prop:7-ols}
$\hat\tau$ 和 $\hat V$ 等于将向量 $(\hat\tau_1, \dots, \hat\tau_n)^\top$ 仅对截距进行 OLS 拟合的系数和方差估计量。
\end{Proposition}

$\hat\tau$ 和 $\hat V$ 的 OLS 解释使得计算极为便捷，只需在 R 中运行 \texttt{lm(tau\_hat \~{} 1)} 即可。\footnote{证明见附录 \cref{sec:appendix-7-ols-proof}}

\section{协变量调整}\label{sec:7-covariate}

虽然在设计阶段我们已经根据协变量进行了配对，但配对可能不完美（$X_{i1} \neq X_{i2}$），有时我们还有配对阶段使用的协变量之外的额外协变量。在这些情况下，我们可以调整协变量以进一步提高估计效率。\footnote{详细推导见附录 \cref{sec:appendix-7-covariate}}

假设每个单元 $(i,j)$ 都有协变量 $X_{ij}$。与前述讨论类似，MPE 中协变量调整有两种一般策略。

\subsection{FRT 中的协变量调整}\label{sec:7-frt-covariate}

从 FRT 角度，我们可以在协变量上拟合结果的 OLS 得到残差 $\hat\varepsilon_{ij}$。由于这些残差在尖锐零假设下是固定数，我们可以将 $\hat\varepsilon_{ij}$ 视为观测结果来构造检验统计量。\citet{Rosenbaum:2002a} 特别倡导这种 MPE 策略。

我们还可以直接使用模型拟合的系数作为检验统计量。

\subsection{回归调整}\label{sec:7-regression-adjustment}

现在我们关注 $\tau$ 的估计。我们可以计算协变量的配对内差异 $\hat\tau_{X,i}$ 及其均值 $\hat\tau_X$，方法与结果相同。可以证明
\[
\mathbb{E}(\hat\tau_{X,i}) = 0, \quad \mathbb{E}(\hat\tau_X) = 0,
\]
和
\[
\text{cov}(\hat\tau_X) = n^{-2} \sum_{i=1}^n \hat\tau_{X,i} \hat\tau_{X,i}^\top.
\]

在实现的 MPE 中，除非所有 $\hat\tau_{X,i}$ 都为零，否则 $\text{cov}(\hat\tau_X)$ 不为零。由于 $(Z_1, \ldots, Z_n)$ 的不幸运抽取，$\hat\tau_X$ 可能显著偏离零。类似于第 \ref{ch:6} 章的讨论，调整协变量均值的不平衡可能会提高估计效率。

类似于第 \ref{ch:6} 章的讨论，我们可以考虑由 $\gamma$ 索引的一类估计量：
\begin{equation}
\hat{\boldsymbol\tau}(\gamma) = \hat{\boldsymbol\tau} - \boldsymbol\gamma^\top \hat{\boldsymbol\tau}_X.
\label{eq:7-adjusted-estimator}
\end{equation}
对于任意固定 $\gamma$，该估计量的均值为零。我们希望选择 $\gamma$ 使 $\hat\tau(\gamma)$ 的方差最小化。

\begin{Theorem}[协方差的无偏估计]\label{th:7-cov-estimate}
$\text{cov}(\hat\tau_X, \hat\tau)$ 的一个无偏估计为
\begin{equation}
\hat\theta = \{n(n-1)\}^{-1} \sum_{i=1}^n (\hat\tau_{X,i} - \hat\tau_X)(\hat\tau_i - \hat\tau).
\label{eq:7-cov-estimator}
\end{equation}
\end{Theorem}

因此，我们可以估计最优系数：
\[
\hat\gamma = \left(n^{-2} \sum_{i=1}^n \hat\tau_{X,i} \hat\tau_{X,i}^\top\right)^{-1} \hat\theta \approx \left(\sum_{i=1}^n (\hat\tau_{X,i} - \hat\tau_X)(\hat\tau_{X,i} - \hat\tau_X)^\top\right)^{-1} \sum_{i=1}^n (\hat\tau_{X,i} - \hat\tau_X)(\hat\tau_i - \hat\tau),
\]
这大约是 $\hat\tau_i$ 对 $\hat\tau_{X,i}$（带截距）进行 OLS 拟合时 $\hat\tau_{X,i}$ 的系数。最终估计量为
\begin{equation}
\hat\tau_{\text{adj}} = \hat\tau(\hat\gamma) = \hat\tau - \hat\gamma^\top \hat\tau_X.
\label{eq:7-final-adjusted}
\end{equation}
由 OLS 的性质，$\hat\tau_{\text{adj}}$ 大约等于将 $\hat\tau_i$ 对 $1$ 和 $\hat\tau_{X,i}$（带截距）进行 OLS 拟合的截距。

$\hat\tau_{\text{adj}}$ 的保守方差估计量为
\[
\hat V_{\text{adj}} = \hat V + \hat\gamma^\top \text{cov}(\hat\tau_X) \hat\gamma - 2\hat\gamma^\top \hat\theta = \hat V - \hat\theta^\top \text{cov}(\hat\tau_X)^{-1} \hat\theta.
\]

\begin{Proposition}[调整估计量的 OLS 解释]\label{prop:7-adjusted-ols}
在 MPE 下，协变量调整估计量 $\hat\tau_{\text{adj}}$ 和相关方差估计量 $\hat V_{\text{adj}}$ 可以方便地近似为将 $\hat\tau_i$ 的向量对 $1$ 的向量和 $\hat\tau_{X,i}$ 的矩阵进行 OLS 拟合的截距和相关方差估计量。
\end{Proposition}

有趣的是，\cref{prop:7-ols} 和 \cref{prop:7-adjusted-ols} 都不需要 EHW（Eicker--Huber--White）校正。因为我们将 MPE 的数据简化为配对内差异，所以不需要像 CRE 下 Lin (2013) 估计量那样对协变量进行中心化。

\section{Darwin 数据的 FRT 分析}\label{sec:7-example-darwin}

这是来自 \citet{Fisher:1935} 的经典例子。它包含 15 对玉米，其中一对接受异花授粉，另一对接受自花授粉，结果是高度。R 包 \texttt{HistData} 提供了原始数据，其中 \texttt{cross} 和 \texttt{self} 分别是异花授粉和自花授粉下的高度，\texttt{diff} 是它们的差值。

整个 MPE 有 $2^{15} = 32768$ 种可能的治疗分配，这是一个可处理的数目。我们枚举所有可能的治疗分配并计算相应的 $\hat\tau$ 和单边精确 p 值：

\begin{verbatim}
> difference = ZeaMays$diff
> n.pairs = length(difference)
> abs.diff = abs(difference)
> t_obs = mean(difference)
> t_ran = sapply(1:2^15,
+ function(x) {
+ sum(MP_enumerate(x, 15)*abs.diff)
+ })/n.pairs
> pvalue = mean(t_ran >= t_obs)
> pvalue
[1] 0.02633667
\end{verbatim}

精确 p 值为 $0.026$，表明在 5\% 显著性水平下拒绝零假设。\cref{fig:7-randomization-dist} 展示了在 $H_{0F}$ 下 $\hat\tau$ 的精确随机化分布。

\begin{figure}[htbp]
\centering
%\includegraphics[width=0.6\textwidth]{images/949ff4c732aefd8922f653d8285e13bbe4d775c3f87a74750290ddd9818622e8.jpg}
\caption{Darwin 数据下 $\hat\tau$ 的精确随机化分布}
\label{fig:7-randomization-dist}
\end{figure}

\section{儿童电视 Workshop 实验}\label{sec:7-example-ctw}

我们重新分析来自 \citet{Ball:1973} 的儿童电视 Workshop 实验数据（也见 \citealp{ImbensRubin:2015} 第 10 章）。它包含 8 对班级，其中一对中的两个班级被随机分配在标准阅读课期间观看《电动教室》节目或不观看。数据包含前测分数（协变量）和后测分数（结果）。

下表总结了对内协变量和结果及其差异：

\begin{center}
\begin{tabular}{c|ccc|ccc}
配对 & \multicolumn{3}{c}{前测} & \multicolumn{3}{c}{后测} \\
 & 班级1 & 班级2 & 差值 & 班级1 & 班级2 & 差值 \\ \hline
1 & 24 & 25 & -1 & 51 & 56 & -5 \\
2 & 42 & 40 & 2 & 64 & 65 & -1 \\
3 & 52 & 52 & 0 & 75 & 72 & 3 \\
4 & 71 & 70 & 1 & 79 & 77 & 2 \\
5 & 29 & 27 & 2 & 54 & 49 & 5 \\
6 & 50 & 50 & 0 & 75 & 78 & -3 \\
7 & 69 & 69 & 0 & 91 & 92 & -1 \\
8 & 36 & 34 & 2 & 64 & 55 & 9
\end{tabular}
\end{center}

基于 Neymanian 推断，我们计算点估计和方差估计。\footnote{完整计算见附录 \cref{sec:appendix-7-ctw-example}}

\section{MPE 与 CRE 的比较}\label{sec:7-mpe-vs-cre}

配对随机实验（MPE）本质上是一种特殊的分层随机实验（SRE），其中每个层恰好包含一个治疗单元和一个对照单元。那么，MPE 与在相同总体上进行完全随机实验（CRE）相比，效率如何？

考虑一个包含 $2n$ 个单元的总体。设 $X_i$ 为配对 $i$ 中两个单元的协变量均值，即 $X_i = (X_{i1} + X_{i2}) / 2$。记配对内协变量差异为 $d_i = X_{i1} - X_{i2}$，其方差为 $\sigma_d^2 = n^{-1}\sum_{i=1}^n d_i^2$。

为比较 MPE 与 CRE，我们需要对潜在结果的协变量依赖进行建模。假设线性结果模型：
\[
Y_{ij}(1) = \alpha + \beta X_{ij} + \varepsilon_{ij}, \quad Y_{ij}(0) = \alpha - \beta X_{ij} + \varepsilon_{ij},
\]
其中 $\varepsilon_{ij}$ 是均值为 $0$、方差为 $\sigma_\varepsilon^2$ 的独立误差项，且与协变量 $X_{ij}$ 无关。在此模型下，单元 $ij$ 的因果效应为 $\tau_{ij} = 2\beta X_{ij}$，平均因果效应为 $\tau = 2\beta \bar{X}$。

\begin{Theorem}[MPE 与 CRE 的效率比较]\label{th:7-mpe-vs-cre}
在上述线性结果模型下，当 $n$ 较大且 $\sigma_d^2 > 0$ 时，MPE 的方差严格小于 CRE 的方差：
\[
\text{var}_\MPE(\hat\tau) < \text{var}_\CRE(\hat\tau).
\]
效率增益来自配对设计消除的配对间变异。
\end{Theorem}

直观上，配对设计利用协变量信息将相似的单元配对，使得配对内两个单元的协变量接近（$|d_i|$ 较小）。由于因果效应与协变量相关，配对内两个单元的因果效应也接近，从而配对内差异 $\hat\tau_i$ 的方差较小。

\cref{th:7-mpe-vs-cre} 表明，\textbf{只要协变量能够预测潜在结果，配对设计就能带来效率提升。} 这一结论与第 \ref{ch:5} 章关于 SRE 与 CRE 比较的结论一致——MPE 是 SRE 的极端形式，同样受益于协变量的使用。

\begin{Proposition}[效率增益的量化]\label{prop:7-efficiency-gain}
设配对内协变量差异的方差为 $\sigma_d^2$，配对数均值的方差为 $\sigma_X^2$。在上述线性模型下，MPE 与 CRE 的方差之比为
\[
\frac{\text{var}_\MPE(\hat\tau)}{\text{var}_\CRE(\hat\tau)} \approx \frac{\sigma_\varepsilon^2}{\sigma_\varepsilon^2 + \beta^2 \sigma_d^2}.
\]
当 $\sigma_d^2 \to 0$（完美配对）时，该比值趋近于 $0$，效率增益最大；当 $\sigma_d^2 \to \infty$（随机配对）时，该比值趋近于 $1$，效率增益消失。
\end{Proposition}

从实践角度，以下因素会影响 MPE 相对于 CRE 的效率增益：
\begin{enumerate}
  \item \textbf{协变量的预测能力}：$\beta$ 越大，协变量对结果越有预测力，效率增益越大；
  \item \textbf{配对质量}：$\sigma_d^2$ 越小，配对越精确，效率增益越大；
  \item \textbf{样本量}：当 $n$ 较小时，有限样本校正可能影响精确比较。
\end{enumerate}

\section{一般配对实验的扩展}\label{sec:7-general-matched}

前述 MPE 的讨论基于一个关键假设：每对恰好包含一个治疗单元和一个对照单元。本节将此设定推广到更一般的情形。

\subsection{每对多个单元}\label{sec:7-general-multi}

有时出于实际考虑，配对内的单元数可能大于 $2$（即每对有多个治疗单元和多个对照单元）。设第 $i$ 对有 $n_{i1}$ 个治疗单元和 $n_{i0}$ 个对照单元，其中 $n_{i1} + n_{i0} = n_i$。

与 \cref{sec:7-neymanian} 类似，配对 $i$ 内的平均因果效应为
\[
\tau_i = n_i^{-1} \sum_{j=1}^{n_i} \{Y_{ij}(1) - Y_{ij}(0)\},
\]
总体平均因果效应为 $\tau = n^{-1}\sum_{i=1}^n n_i \tau_i$。

配对 $i$ 内的差值均值估计量为
\[
\hat\tau_i = \bar{Y}_{i1}(\text{obs}) - \bar{Y}_{i0}(\text{obs}),
\]
其中 $\bar{Y}_{i1}(\text{obs})$ 和 $\bar{Y}_{i0}(\text{obs})$ 分别是治疗组和对照组观测结果的均值。

加权估计量为
\[
\hat\tau = \sum_{i=1}^n \frac{n_i}{n} \hat\tau_i.
\]

方差估计量需要考虑配对内随机化。类似 \cref{sec:7-neymanian}，我们可以证明方差估计量的保守性。

\subsection{不完美配对}\label{sec:7-imperfect-pairing}

前述讨论假设配对完美，即同一对内两个单元的协变量完全相同。但实际中，由于协变量的连续性或维度较高，我们可能无法找到精确匹配的单元。此时，每对内的协变量存在差异 $(X_{i1} \neq X_{i2})$。

不完美配对下的 MPE 仍然有效，因为随机化仍然保证了比较的基础。但估计精度会有所下降。直观上，配对越不完美（$|d_i|$ 越大），MPE 越接近 CRE，效率增益越小。

\begin{Remark}
当配对不完美时，第 \ref{sec:7-covariate} 节的协变量调整策略尤为重要。通过调整配对内协变量差异，我们可以恢复部分效率损失。
\end{Remark}

\subsection{层内配对}\label{sec:7-stratified-pairing}

另一种扩展是先按协变量分层，再在每层内进行配对。这种两层设计结合了 SRE（减少层间变异）和 MPE（减少层内变异）的优点。

设 $K$ 个层，第 $k$ 层有 $n_{[k]}$ 个单元。先按 $X$ 将单元分层，再在每层内将单元配对。最终的估计量是层别加权平均：
\[
\hat\tau = \sum_{k=1}^K \pi_{[k]} \hat\tau_{[k]},
\]
其中 $\hat\tau_{[k]}$ 是第 $k$ 层内的配对估计量。

这种设计在协变量高度预测结果时尤其有效，因为可以在关键的协变量层内精确匹配单元。

\section{本章小结}\label{sec:7-summary}

本章介绍了配对随机实验（MPE），主要内容如下：

\begin{enumerate}
  \item \textbf{MPE 的动机}：MPE 是 SRE 的最极端形式，当每层只有一个治疗单元和一个对照单元时，配对设计可以最大程度消除层间变异。
  \item \textbf{MPE 定义（\cref{def:7-mpe}）}：治疗指示变量 $Z_i \stackrel{\text{IID}}{\sim} \text{Bernoulli}(1/2)$。
  \item \textbf{FRT}：可以使用多种检验统计量，包括配对 $t$ 统计量（\cref{eq:7-t-pair}）、Wilcoxon 符号秩统计量（\cref{eq:7-wilcoxon}）、Kolmogorov--Smirnov 型统计量（\cref{eq:7-ks-statistic}）和符号统计量（\cref{eq:7-sign-stat}）。
  \item \textbf{Neymanian 推断}：方差估计量 $\hat V$（\cref{eq:7-variance-estimator}）是保守的（\cref{th:7-variance-conservative}）；点估计和方差估计都可以通过 OLS 方便地得到（\cref{prop:7-ols}）。
  \item \textbf{协变量调整}：调整估计量 $\hat\tau_{\text{adj}}$（\cref{eq:7-final-adjusted}）和方差估计量 $\hat V_{\text{adj}}$ 可以通过 OLS 方便地计算（\cref{prop:7-adjusted-ols}）。
  \item \textbf{MPE 与 CRE 的比较}：只要协变量能预测潜在结果，MPE 的方差严格小于 CRE 的方差（\cref{th:7-mpe-vs-cre}）。
  \item \textbf{一般配对实验}：MPE 可以推广到每对多个单元、不完美配对和层内配对等情形。
\end{enumerate}

\section{习题}\label{sec:chapter7-exercises}

\begin{Exercise}{MPE 的方差}\label{exr:7-1}
证明在 MPE 下，$\hat\tau$ 的方差为
\[
\text{var}(\hat\tau) = n^{-2} \sum_{i=1}^n \text{var}(\hat\tau_i).
\]
\end{Exercise}

\begin{Exercise}{配对 $t$ 统计量的渐近分布}\label{exr:7-2}
证明 \cref{eq:7-t-asymptotic}：在 $H_{0F}$ 下，
\[
\frac{\hat\tau}{\sqrt{n^{-2} \sum_{i=1}^n \hat\tau_i^2}} \to \mathrm{N}(0, 1)
\]
依分布收敛。
\end{Exercise}

\begin{Exercise}{Wilcoxon 符号秩统计量的矩}\label{exr:7-3}
证明 \cref{ex:7-wilcoxon} 中 $\mathbb{E}(W)$ 和 $\text{var}(W)$ 的公式。
\end{Exercise}

\begin{Exercise}{McNemar 统计量的精确分布}\label{exr:7-4}
证明在 $H_{0F}$ 下，$m_{10} \sim \text{Binomial}(m_{10} + m_{01}, 1/2)$。
\end{Exercise}

\begin{Exercise}{方差估计量的偏}\label{exr:7-5}
证明 \cref{th:7-variance-conservative}：在 MPE 下，
\[
\mathbb{E}(\hat V) - \text{var}(\hat\tau) = \{n(n-1)\}^{-1} \sum_{i=1}^n (\tau_i - \tau)^2 \geq 0.
\]
\end{Exercise}

\begin{Exercise}{OLS 解释}\label{exr:7-6}
证明 \cref{prop:7-ols}：$\hat\tau$ 和 $\hat V$ 等于将 $(\hat\tau_1, \dots, \hat\tau_n)^\top$ 仅对截距进行 OLS 拟合的系数和方差估计量。
\end{Exercise}

\begin{Exercise}{协方差估计量}\label{exr:7-7}
证明 \cref{th:7-cov-estimate}：$\text{cov}(\hat\tau_X, \hat\tau)$ 的无偏估计为
\[
\hat\theta = \{n(n-1)\}^{-1} \sum_{i=1}^n (\hat\tau_{X,i} - \hat\tau_X)(\hat\tau_i - \hat\tau).
\]
\end{Exercise}

\begin{Exercise}{调整估计量的 OLS 解释}\label{exr:7-8}
证明 \cref{prop:7-adjusted-ols}：协变量调整估计量 $\hat\tau_{\text{adj}}$ 可以近似等于将 $\hat\tau_i$ 对 $1$ 和 $\hat\tau_{X,i}$（带截距）进行 OLS 拟合的截距。
\end{Exercise}

\begin{Exercise}{Darwin 数据的 Neymanian 分析}\label{exr:7-9}
使用 Neymanian 方法重新分析 Darwin 数据。报告点估计、方差估计和 95\% 置信区间。
\end{Exercise}

\begin{Exercise}{儿童电视 Workshop 实验的 FRT 分析}\label{exr:7-10}
使用 FRT 重新分析儿童电视 Workshop 实验数据，使用配对 $t$ 统计量作为检验统计量。
\end{Exercise}


